\section{Sibling signing and sibling graph separation}

The sibling certificate pair addresses a nearby signed carrier variant. The point is not to invalidate the sibling, but to show that it is genuinely distinct from the canonical signing and from the canonical graph. This sibling is retained because it tests whether a small carrier-cycle change can preserve headline scale while changing the actual global graph.

\subsection{Non-switching certificate}

The sibling signing is defined as the canonical signing plus a delta supported on six slot edges:
\[
(5,10),(5,14),(10,14),(9,11),(9,12),(11,12).
\]
Equivalently, these six signs are toggled and all other slot-edge signs are left unchanged. The support consists of two triangle supports:
\[
(5,10,14),\qquad (9,11,12).
\]

Let \(B\) be the vertex-edge coboundary matrix of the slot graph over \(\mathbb F_2\). The columns of \(B\) are indexed by the 15 slot vertices and the rows are indexed by the 30 slot edges. For an edge \(e=\{u,v\}\), the row of \(B\) has ones in the \(u\)- and \(v\)-columns and zeros elsewhere. A switching function is a vector \(g\in\mathbb F_2^{15}\), and the signing difference is a vector \(\delta\in\mathbb F_2^{30}\). With this convention, switching equivalence is tested by the equation
\[
Bg=\delta.
\]

The certificate reports that this equation is not solvable:
\[
\operatorname{rank}(B)=14,
\qquad
\operatorname{rank}([B|\delta])=15.
\]
Since the augmented rank exceeds the coefficient rank, the sibling delta is not a coboundary. Therefore the sibling signing is not switching-equivalent to the canonical signing.

\subsection{Full graph separation}

The sibling full graph separation certificate builds the sibling graph from its signing table and compares exact graph invariants.

\begin{center}
\begin{tabular}{lr}
\toprule
Quantity & Value \\
\midrule
Canonical edges & 3600 \\
Sibling edges & 3600 \\
Intersection edges & 3240 \\
Canonical-only edges & 360 \\
Sibling-only edges & 360 \\
Symmetric difference & 720 \\
\bottomrule
\end{tabular}
\end{center}

The two graphs share diameter and radius:
\[
\operatorname{diam}=8,\qquad \operatorname{rad}=6.
\]
However, their eccentricity counts and distance distributions differ. Since distance distribution is a graph invariant, this proves non-isomorphism for the recorded sibling graph.

\begin{center}
\begin{tabular}{lrr}
\toprule
Eccentricity & Canonical count & Sibling count \\
\midrule
6 & 349 & 327 \\
7 & 541 & 525 \\
8 & 10 & 48 \\
\bottomrule
\end{tabular}
\end{center}

\begin{center}
\begin{tabular}{lrr}
\toprule
Distance & Canonical pairs & Sibling pairs \\
\midrule
1 & 3600 & 3600 \\
2 & 17710 & 17700 \\
3 & 60345 & 60642 \\
4 & 131446 & 130810 \\
5 & 143177 & 142395 \\
6 & 45600 & 46225 \\
7 & 2667 & 3144 \\
8 & 5 & 34 \\
\bottomrule
\end{tabular}
\end{center}

